\documentclass[12pt,a4paper]{report}
\usepackage[utf8]{inputenc}
\usepackage[T1]{fontenc}
\usepackage[french]{babel}
\usepackage{geometry}
\geometry{margin=2.5cm}

\begin{document}

\begin{titlepage}
    \centering
    \vspace*{4cm}
    {\Huge Kernax}\\[0.5cm]
    {\large Programmer's Guide}\\[2cm]
    {\large Ziad Iziz}\\[0.5cm]
    {\large INFOB318 --- UNamur}\\[0.5cm]
    {\large 2026}
\end{titlepage}

\tableofcontents
\newpage

\chapter{Contexte et objectifs techniques}
\section{Pourquoi JAX ?}
Le choix de JAX plutot que pytorch a été motivé par l'optimisation des paramètres des kernels.
JAX permet d'appliquer value_and_grad direct sur le log vraisemblance, ce qui nous simplifie l'implem.
\section{Limites connues}
Mon SVM n'est pas instancié comme pytree, erreur, ce qui veut dire qu'on limite les possibilités de jit.
\chapter{Environnement de développement}
\section{Stack}
\section{Générer la documentation}

\chapter{Architecture globale}
\section{Vue d'ensemble}
\section{Le pattern pytree}

\chapter{Description des modules}
\section{GP.py}
\section{SVM.py}
\section{KRR.py}

\chapter{Référence de l'API}

\chapter{Comportement au runtime}
\section{Flux de fit()}
\section{Flux de predict()}

\chapter{Pistes d'amélioration}

\end{document}